\section{四维证明的技术骨架}

四维篇的证明是一个反证：假设存在 Hausdorff 维数小于 $4$ 的挂谷集，构造一个带正“时间代价”的加权问题，再从这个问题中逼出一个自相矛盾的“改进”。本章按证明自身的模块顺序重建骨架：归约（第 2 章）、比较与延拓（第 3 章）、极值临界路径（第 4 章）、三个窄度区间（第 6、8--10 章）、收口与推论（第 10、11 章）。术语第一次出现时附原文，便于对照稿件。

\subsection{归约：把任意集合问题变成加权线族问题}

归约的定理陈述可以一字不改地引用\upcite{openai4d2026}：

\begin{theorem}[Reduction, 四维稿件 Theorem 2.17，转述]
若 $K\subset\mathbb{R}^4$ 不含任何可测性或紧性假设，包含每个方向的单位线段，且 $\dim_{\mathrm{H}}K<4$，则存在 $d<1$、$S,\eta>0$，以及一个属于容量类 $\mathcal{C}(d)$ 的温和齐次精确问题 $E$（版本 2），使得 $H(E)>0$。
\end{theorem}

这里连续出现了三个机制性概念，逐个解释。

\textbf{加权模型}。把“集合含每方向线段”翻译成概率语言：取一族轨迹（trajectory）$(t,\; y_0+tV,\; w_0+tW)$（其中 $y\in\mathbb{R}^2$ 是基础坐标，$w$ 是隐藏坐标），每条轨迹服从某个先验测度，时间独立均匀；再用一个入射选择（incidence selection）标记哪些“轨迹--时间”对被保留。这样处理后，问题从“集合的覆盖数”转成“加权入射测度的分布”，密度、角度集中度这类量都可以用指数记账。归约的第一步用 Hausdorff 覆盖的外测度可数次可加性，挑出一个斜率正测度、截距有界的子族——这一步不需要可测性，是“任意集合”能进入框架的关键。

\textbf{观察与密度}。在深度 $r$ 上做分箱观察：保留精确的基础坐标，砍掉隐藏坐标（记为 $[w]_{a_r/B}$），得到一组嵌套的观测 $p_r$。入射密度沿深度按 $n(r)=N_0^{-f(r)}$ 退化。反证假设 $\dim_{\mathrm{H}}K<4$ 体现在密度退化速度上：存在 $d<1$，使 $f(L)-f(r)\le d(L-r)$，即深度每推进一格，密度损失不到一个完整幂（容量类 $\mathcal{C}(d)$ 的第一条）。第二条是角度非集中：隐藏观察细化时，速度方向的分布不许过度集中。

\textbf{图卡、覆盖与地平线}。图卡（chart）是证明的局部对象：一个移动的空间盒、一个时间区间 $I$，加上一个二次多项式测试 $P$，要求 $P$ 在足够多的被选轨迹上保持小（厚度受控）。图卡自身带质量要求，用“选定的时间轴对”度量，防止“只拟合了可忽略的例外族”被当成进展。区间长度写成 $N_0^{-h+o(1)}$ 时，$h$ 叫时间代价（time cost）；一个图卡系统的覆盖量是其选定入射质量，质量 $N_0^{-o(1)}$ 叫实质覆盖。于是：

\begin{equation}
H(E)\;=\;\inf\{\, h \;:\; \text{存在带实质覆盖的真终端图卡系统，区间长度} = N_0^{-h+o(1)} \,\},
\end{equation}

即地平线（horizon）——在“拟合到最终厚度”这件事上，时间局部化至少要付的代价。归约的结论是 $H(E)>0$：反例逼出了一笔正代价。（地平线的形式定义、图卡与真图卡的完整条件、以及归约定理的形态，转录于附录 C.1，供逐条核对。）

\begin{keybox}[归约的要点]
反证假设（维数亏空）被换算成一条关于指数的正下界（地平线为正）。此后全部工作变成：证明 $H(E)=0$，或者证明 $H(E)>0$ 会导致矛盾。整个四维证明的剩余部分都在执行后一种路线。
\end{keybox}

\subsection{比较原理：跨尺度把局部描述接起来}

有了图卡，下一步是把不同尺度上的局部描述拼接起来。第 3 章的比较原理（comparison of algebraic sheets）处理的正是拼接的合法性：多项式测试的局部图像给出“标量叶片”（scalar sheets），即多项式方程的解析分支（精确定义与 Lemma 3.2、Lemma 3.3 的定量陈述见附录 C.1）。密度约束与保留下来的低界质量给出一个关键定量事实：在精度 $\delta$ 上，每个局部描述只需一个大小至多 $\delta^{-d+o(1)}$ 的列表——由于 $d<1$，列表规模小于任何完整幂。

比较原理的机制值得复述：两个这样的列表不可能“频繁地取值一致、同时保持换算后的导数数据分离”；否则，交替抽取新条件下的叶片与时间，会把持续的导数失配拼成一个多项式端点映射，其像的体积超过密度约束允许的上限。成功路径保持为非归一化子测度，因此体积论证保留了对先前选择的依赖——这个技术细节保证比较不会悄悄放宽假设。

\begin{evbox}[证据：d<1 与比较原理的耦合]
容量类 $\mathcal{C}(d)$ 要求 $d<1$，稿件给出的原因正是比较原理：列表规模 $\delta^{-d+o(1)}$ 只有在 $d<1$ 时才能被“体积超界”论证压住\upcite{openai4d2026}。这类“一个常数被证明机制的结构所选定”的安排，是手稿设计感的典型体现。
\end{evbox}

\subsection{极值选择与临界路径}

第 4 章把问题摆到一个近乎极值的加权问题上：沿着嵌套图卡行进，相对密度与时间剖面会线性化，出现正整数的时间速率 $k>0$；再做一次极值选择，控制“较小的空间轴能比时间尺度细多少”，把最优窄度（optimized narrowness）记为 $\ell$。$\ell$ 大体上衡量管族压扁的方向——$\ell=0$ 是各向同性，$\ell$ 有限是有限窄度，$\ell=\infty$ 是无界窄度。

“改进”（improvement）的定义由此而来：要么在保持覆盖的情况下把时间代价严格推低，要么在 $\ell$ 有限时于固定时间上换取额外的窄度增益。极值选择保证：任何一处这样的改进，回到原问题后都违反极值性。于是证明的战术目标变成了“在每一个实质性残余上制造一次改进”。

从局部候选到“实际改进”的转换是第 4 章的另一半工作，收在命题 Proposition 4.11 与引理 Lemma 4.12：先用有限贪心赋值把候选情景压实成有实质覆盖的单一情景，再用小列表比较把候选测试组装成改进的图卡测试；最后 Lemma 4.12 把有限列表从分析用子测度回填到原测度（误差由基网格上的全变差控制）。这一步是整套证明里最容易“看起来像行政细节、实际是逻辑要冲”的地方——没有它，局部拟合无法升级为对原问题成立的改进。

\subsection{三个窄度区间}

\begin{table}[t]
\centering
\caption{三个窄度区间：假设、机制与产出}
\label{tab:regimes}
\footnotesize
\begin{tabularx}{\textwidth}{@{}L{1.9cm}L{3.4cm}Y@{}}
\toprule
\rowcolor{lightbg}\textbf{区间} & \textbf{假设} & \textbf{机制与产出}（依据稿件第 6、8--10 章） \\
\midrule
$\ell=0$ 各向同性 & 分离速度足够多 & 沿单线的条件标量估计给出多项式拟合；多元 Kakeya 与插值合并；换线迫使近似守恒，产出势函数；宽带方向的圆锥障碍由隐藏坐标编码三次、四次项修回二次测试（次数恢复，Proposition 6.8） \\
$0<\ell<\infty$ 有限窄度 & 投影后部分区间足够宽 & 宽投影区间直接由比较原理改进；否则平面刚性把轨迹压入水平模型 $Y'=-Z+tv$，自然盒型为水平尺度 $H$、横向 $H^{2}$；临界速率 $2k=1$ 时三条刷新线段填满三维基础体积，偏离临界速率时回归宽区间，$2k<1$ 时先做自举 \\
$\ell=\infty$ 无界窄度 & 窄度趋于无穷 & 两个嵌套图卡标签把轨迹在全归一化时间区间上局部化；归一化造成的密度指数大漂移在相对比值中抵消；最终标量投影产出时间代价趋于零的图卡，与正地平线矛盾 \\
\bottomrule
\end{tabularx}
\end{table}

三个区间共享一条供给线：各向同性情形的标量版本（单空间坐标）先被证明“零地平线”，其副产品是标量投影定理（Theorem 8.2），供另外两个区间做降维使用\upcite{openai4d2026}。这种“先在最难区间把一维辅助模型做出来，再让其他区间站上去”的安排，使三区间不是三条平行线，而是一棵树。

表 \ref{tab:regimes} 中两处机制值得展开一句。其一，圆锥障碍与次数恢复：各向同性情形产出的势函数初始次数可能高于 $2$，而“宽带方向落在圆锥附近”是唯一无法直接排除的障碍形态；稿件识别出残余的三次、四次项，用原本闲置的隐藏坐标把它们编码进二次修正，从而恢复图卡所需的次数。其二，水平模型与临界速率：$Y'=-Z+tv$ 这一形式的水平盒为 $H\times H^{2}$，其“非交换”的水平方向迫使标量叶片近似仿射，从而让 $2k=1$ 的三腿切换恰好把三维基础体积填满。

\subsection{收口、几何推论与三维范式的对照}

合成段（第 10 章）把三个区间的产出拼回原问题：每个区间先在加权族的实质残余上产出局部候选，再经贪心覆盖与小列表比较升级为可回填到原问题的改进图卡系统，与极值性矛盾。归纳法与 $H(E)>0$ 的阳性下界共同给出 $\dim_{\mathrm{H}}K\ge 4$；上界是空间的平凡上界，定理闭合\upcite{openai4d2026}。

第 11 章把主定理摊成一组推论：有界挂谷集同时具有满的下、上 Minkowski 维数与 packing 维数；对 $d\ge 4$，经四维子空间的正交投影可得 $\dim_{\mathrm{H}}E\ge 4$（$d=5$ 也包含在内，但离 $5$ 的完整猜想还差得远）；任意方向集版本 $\dim_{\mathrm{H}}E\ge 1+\dim_{\mathrm{H}}D$ 经 Keleti--M\'ath\'e 的转移得到；常曲率流形上的 Nikodym 集与固定平移不变相下的弯曲挂谷族，经 Gao--Liu--Xi 的转移定理得到。这些推论都附带各自的假设条件，稿件据此逐条列出\upcite{openai4d2026,keletimathe2023,gaoliuxi2025}。

\begin{table}[t]
\centering
\caption{四维稿件与三维范式在对象层面的对照}
\label{tab:3d4d}
\footnotesize
\begin{tabularx}{\textwidth}{@{}L{2.8cm}L{4.6cm}Y@{}}
\toprule
\rowcolor{lightbg}\textbf{维度} & \textbf{结构分类} & \textbf{证明载体} \\
\midrule
三维（Wang--Zahl 传统） & 粘性、平面性、颗粒性（近极端管族的三种形态） & 棱柱与阴影加权、多尺度归纳、凸聚类；关键定理在四维失效（二次曲面可容纳三维线族） \\
四维（OpenAI 手稿） & 按窄度参数 $\ell$ 划三区间，取代形态分类 & 二次图卡与质量约束、地平线极值反证、标量投影降维、水平模型 \\
\bottomrule
\end{tabularx}
\end{table}

表 \ref{tab:3d4d} 是本文对两代方法差异的概括（本文判断，\cfM）：三维的范式是“先分类极端形态，再逐一击破”，四维稿件的范式是“不分类形态，改为对局部拟合的时间代价做极值分析”——后者对窄度参数 $\ell$ 的连续性处理，正是为了绕开四维中比三维复杂得多、可能穷举不尽的聚集形态。

\begin{keybox}[可记忆主句]
四维证明的骨架一句话概括：反例逼出一条正地平线，三区间各自证明能在实质残余上制造一次改进，而改进与极值性不相容——于是反例不能存在。
\end{keybox}

\subsection{三维姊妹篇的骨架（作为对照）}

三维篇的目标是把阴影密度幂补到 $\lambda^{3}$。其机制与四维篇同源但更紧凑：先建立一个对密度敏感的临界指数 $h(p,z)$，把“极大定理成立”等价为“临界指数为零”；再假设临界指数为正，构造取等序列（equality configurations），用局部质量界与多元 Kakeya 逼出“短块上方向近共面”，用平面 Furstenberg 估计与控制支撑代价的板条 Nikodym 估计把“时间亏空”压成零；随后证明取等配置在跨尺度上必然收敛到一条规范时间剖面 $F(s)=\beta(s-\tau)_{+}$，并在同一指标的两个事件（双帧实验）之间建立投影约束，与熵需求矛盾\upcite{openai3d2026}。这条链条中，“同一根管子取两个时刻”的双帧实验、以及对钉住（pinning）估计的改造，与熵账本一起构成了第 7 章谱系分析的主要对象。

\begin{keybox}[可记忆主句]
三维篇与四维篇共享同一套记账语言（时间代价、覆盖、极值改进），差别在于三维篇把矛盾压在“双帧投影约束”上，四维篇把矛盾摊在“三区间改进”上。
\end{keybox}